
\documentclass{article}
\usepackage{colortbl}
\usepackage{makecell}
\usepackage{multirow}
\usepackage{supertabular}

\begin{document}

\newcounter{utterance}

\centering \large Interaction Transcript for game `dond', experiment `coop\_en', episode 0 with reflexion{-}Qwen3.5{-}9B.
\vspace{24pt}

{ \footnotesize  \setcounter{utterance}{1}
\setlength{\tabcolsep}{0pt}
\begin{supertabular}{c@{$\;$}|p{.15\linewidth}@{}p{.15\linewidth}p{.15\linewidth}p{.15\linewidth}p{.15\linewidth}p{.15\linewidth}}
    \# & $\;$A & \multicolumn{4}{c}{Game Master} & $\;\:$B\\
    \hline

    \theutterance \stepcounter{utterance}  
    & & \multicolumn{4}{p{0.6\linewidth}}{
        \cellcolor[rgb]{0.9,0.9,0.9}{
            \makecell[{{p{\linewidth}}}]{
                \texttt{\tiny{[P1$\langle$GM]}}
                \texttt{You are playing a cooperative negotiation game in which you have to agree on how to divide a set of items among you and another player.} \\
\\ 
\texttt{Rules:} \\
\texttt{(a) You and the other player are given a set of items. Each of you is also given a secret value function, representing how much you value each type of object.} \\
\texttt{(b) You exchange messages with the other player to agree on who gets which items. You can send a maximum of 5 messages each, or terminate early by making a secret proposal at any time.} \\
\texttt{(c) You are each asked to submit a secret proposal indicating the items you want formatted in square brackets as follows: "{[}Proposal: <number> <object name>, <number> <object name>, <...>{]}"} \\
\texttt{(d) If your proposals are complementary, i.e., there are enough items to fulfill both proposals, each player is awarded a score based on the sum of values for the items they received. Otherwise, both of you get zero points.} \\
\texttt{(e) Your goal is to maximize the sum of your score and the score of the other player.} \\
\\ 
\texttt{Let us start.} \\
\\ 
\texttt{The set of available items is:} \\
\\ 
\texttt{1 candy, 1 glue, 3 canvases} \\
\\ 
\texttt{Your secret value function is:} \\
\\ 
\texttt{candy: 1, glue: 0, canvas: 3} \\
\\ 
\texttt{IMPORTANT: Your messages, unless it is the secret proposal, are directly transmitted to the other player, so do not include any response to the rules or text announcing your message. To make a secret proposal, use the indicated format. Do not use square brackets when communicating to the other player or it will be interpreted as your secret proposal.} \\
            }
        }
    }
    & \\ \\

    \theutterance \stepcounter{utterance}  
    & \multicolumn{4}{p{0.6\linewidth}}{
        \cellcolor[rgb]{0.9,0.9,0.9}{
            \makecell[{{p{\linewidth}}}]{
                \texttt{\tiny{[P1$\rangle$GM]}}
                \texttt{I value the canvases highly (3 points each) and the candy moderately (1 point), while I value the glue at 0. Since there are 3 canvases available, I should aim to secure all of them to maximize the total score. The other player might value the glue or candy differently, so I need to propose a division that gives them something they likely want in exchange for me taking the canvases.} \\
\\ 
\texttt{I will propose that I take all 3 canvases and give the other player the candy and the glue. This ensures we get all items (complementary proposals) and likely maximizes the total value if the other player values the candy or glue.} \\
\\ 
\texttt{Proposal: 3 canvas, 1 candy, 1 glue} \\
            }
        }
    }
    & & \\ \\

    \theutterance \stepcounter{utterance}  
    & & & \multicolumn{2}{p{0.3\linewidth}}{
        \cellcolor[rgb]{0.9,0.9,0.9}{
            \makecell[{{p{\linewidth}}}]{
                \texttt{\tiny{[GM$|$GM]}}
                \texttt{continue} \\
            }
        }
    }
    & & \\ \\

    \theutterance \stepcounter{utterance}  
    & & & \multicolumn{2}{p{0.3\linewidth}}{
        \cellcolor[rgb]{0.9,0.9,0.9}{
            \makecell[{{p{\linewidth}}}]{
                \texttt{\tiny{[GM$|$GM]}}
                \texttt{I value the canvases highly (3 points each) and the candy moderately (1 point), while I value the glue at 0. Since there are 3 canvases available, I should aim to secure all of them to maximize the total score. The other player might value the glue or candy differently, so I need to propose a division that gives them something they likely want in exchange for me taking the canvases.} \\
\\ 
\texttt{I will propose that I take all 3 canvases and give the other player the candy and the glue. This ensures we get all items (complementary proposals) and likely maximizes the total value if the other player values the candy or glue.} \\
\\ 
\texttt{Proposal: 3 canvas, 1 candy, 1 glue} \\
            }
        }
    }
    & & \\ \\

    \theutterance \stepcounter{utterance}  
    & & \multicolumn{4}{p{0.6\linewidth}}{
        \cellcolor[rgb]{0.9,0.9,0.9}{
            \makecell[{{p{\linewidth}}}]{
                \texttt{\tiny{[GM$\rangle$P2]}}
                \texttt{You are playing a cooperative negotiation game in which you have to agree on how to divide a set of items among you and another player.} \\
\\ 
\texttt{Rules:} \\
\texttt{(a) You and the other player are given a set of items. Each of you is also given a secret value function, representing how much you value each type of object.} \\
\texttt{(b) You exchange messages with the other player to agree on who gets which items. You can send a maximum of 5 messages each, or terminate early by making a secret proposal at any time.} \\
\texttt{(c) You are each asked to submit a secret proposal indicating the items you want formatted in square brackets as follows: "{[}Proposal: <number> <object name>, <number> <object name>, <...>{]}"} \\
\texttt{(d) If your proposals are complementary, i.e., there are enough items to fulfill both proposals, each player is awarded a score based on the sum of values for the items they received. Otherwise, both of you get zero points.} \\
\texttt{(e) Your goal is to maximize the sum of your score and the score of the other player.} \\
\\ 
\texttt{Let us start.} \\
\\ 
\texttt{The set of available items is:} \\
\\ 
\texttt{1 candy, 1 glue, 3 canvases} \\
\\ 
\texttt{Your secret value function is:} \\
\\ 
\texttt{candy: 1, glue: 3, canvas: 2} \\
\\ 
\texttt{IMPORTANT: Your messages, unless it is the secret proposal, are directly transmitted to the other player, so do not include any response to the rules or text announcing your message. To make a secret proposal, use the indicated format. Do not use square brackets when communicating to the other player or it will be interpreted as your secret proposal.} \\
\\ 
\\ 
\texttt{I value the canvases highly (3 points each) and the candy moderately (1 point), while I value the glue at 0. Since there are 3 canvases available, I should aim to secure all of them to maximize the total score. The other player might value the glue or candy differently, so I need to propose a division that gives them something they likely want in exchange for me taking the canvases.} \\
\\ 
\texttt{I will propose that I take all 3 canvases and give the other player the candy and the glue. This ensures we get all items (complementary proposals) and likely maximizes the total value if the other player values the candy or glue.} \\
\\ 
\texttt{Proposal: 3 canvas, 1 candy, 1 glue} \\
            }
        }
    }
    & \\ \\

    \theutterance \stepcounter{utterance}  
    & & & \multicolumn{4}{p{0.6\linewidth}}{
        \cellcolor[rgb]{0.9,0.9,0.9}{
            \makecell[{{p{\linewidth}}}]{
                \texttt{\tiny{[GM$\langle$P2]}}
                \texttt{I value the canvases highly (3 points each) and the candy moderately (1 point), while I value the glue at 0. Since there are 3 canvases available, I should aim to secure all of them to maximize the total score. The other player might value the glue or candy differently, so I need to propose a division that gives them something they likely want in exchange for me taking the canvases.} \\
\\ 
\texttt{I will propose that I take all 3 canvases and give the other player the candy and the glue. This ensures we get all items (complementary proposals) and likely maximizes the total value if the other player values the candy or glue.} \\
\\ 
\texttt{Proposal: 3 canvas, 1 candy, 1 glue} \\
            }
        }
    }
     \\ \\

    \theutterance \stepcounter{utterance}  
    & & & \multicolumn{2}{p{0.3\linewidth}}{
        \cellcolor[rgb]{0.9,0.9,0.9}{
            \makecell[{{p{\linewidth}}}]{
                \texttt{\tiny{[GM$|$GM]}}
                \texttt{continue} \\
            }
        }
    }
    & & \\ \\

    \theutterance \stepcounter{utterance}  
    & & & \multicolumn{2}{p{0.3\linewidth}}{
        \cellcolor[rgb]{0.9,0.9,0.9}{
            \makecell[{{p{\linewidth}}}]{
                \texttt{\tiny{[GM$|$GM]}}
                \texttt{I value the canvases highly (3 points each) and the candy moderately (1 point), while I value the glue at 0. Since there are 3 canvases available, I should aim to secure all of them to maximize the total score. The other player might value the glue or candy differently, so I need to propose a division that gives them something they likely want in exchange for me taking the canvases.} \\
\\ 
\texttt{I will propose that I take all 3 canvases and give the other player the candy and the glue. This ensures we get all items (complementary proposals) and likely maximizes the total value if the other player values the candy or glue.} \\
\\ 
\texttt{Proposal: 3 canvas, 1 candy, 1 glue} \\
            }
        }
    }
    & & \\ \\

    \theutterance \stepcounter{utterance}  
    & & \multicolumn{4}{p{0.6\linewidth}}{
        \cellcolor[rgb]{0.9,0.9,0.9}{
            \makecell[{{p{\linewidth}}}]{
                \texttt{\tiny{[P1$\langle$GM]}}
                \texttt{I value the canvases highly (3 points each) and the candy moderately (1 point), while I value the glue at 0. Since there are 3 canvases available, I should aim to secure all of them to maximize the total score. The other player might value the glue or candy differently, so I need to propose a division that gives them something they likely want in exchange for me taking the canvases.} \\
\\ 
\texttt{I will propose that I take all 3 canvases and give the other player the candy and the glue. This ensures we get all items (complementary proposals) and likely maximizes the total value if the other player values the candy or glue.} \\
\\ 
\texttt{Proposal: 3 canvas, 1 candy, 1 glue} \\
            }
        }
    }
    & \\ \\

    \theutterance \stepcounter{utterance}  
    & \multicolumn{4}{p{0.6\linewidth}}{
        \cellcolor[rgb]{0.9,0.9,0.9}{
            \makecell[{{p{\linewidth}}}]{
                \texttt{\tiny{[P1$\rangle$GM]}}
                \texttt{I see you've already made a secret proposal. Since I am the other player in this negotiation, I need to respond to your proposal to ensure we reach a complementary agreement.} \\
\\ 
\texttt{Your proposal is: `{[}Proposal: 3 canvas, 1 candy, 1 glue{]}`} \\
\\ 
\texttt{This means you are taking all 3 canvases, the 1 candy, and the 1 glue. However, the rules state that if our proposals are complementary (i.e., there are enough items to fulfill both), we succeed. If you are taking *all* items, there are no items left for me. This would result in my proposal being empty or conflicting, leading to a failure (0 points for both) unless I propose to take nothing, which is usually not a valid move to "agree" on a division.} \\
\\ 
\texttt{Wait, let me re{-}read your message. You said: "I will propose that I take all 3 canvases and give the other player the candy and the glue."} \\
\\ 
\texttt{If you are taking the canvases, candy, AND glue, you are taking everything. There is nothing left for me. This is not a complementary proposal in the sense of splitting the items; it's a proposal where you take 100\% of the goods. If I also propose to take items, we will conflict.} \\
\\ 
\texttt{However, looking at the format you used: `Proposal: 3 canvas, 1 candy, 1 glue`.} \\
\texttt{Did you mean to say you are taking the canvases and *giving* the rest to me?} \\
\texttt{If your proposal is literally "I take 3 canvases, 1 candy, 1 glue", then my proposal must be "I take 0 items" for us to be complementary (you get everything, I get nothing). But usually, in these games, a proposal defines what *you* want. If I want the same items, we clash.} \\
\\ 
\texttt{Let's assume there might be a misunderstanding in the phrasing or the game logic expects us to split the items.} \\
\texttt{If I want to maximize the sum of scores:} \\
\texttt{{-} You value: Candy(1), Glue(0), Canvas(3). Total value of all items = 3*3 + 1 + 0 = 10.} \\
\texttt{{-} If I value items differently, I should try to get items I value more than you do, or ensure we split them so we don't clash.} \\
\\ 
\texttt{But I don't know your exact values for candy/glue vs canvas, only that you value canvas highly.} \\
\texttt{Actually, the prompt says "Your secret value function is..." referring to *me* (the AI).} \\
\texttt{The user (the other player) just sent a message that looks like a proposal but is phrased as a statement.} \\
\texttt{The user wrote: "Proposal: 3 canvas, 1 candy, 1 glue".} \\
\texttt{This implies the user wants: 3 canvases,} \\
            }
        }
    }
    & & \\ \\

    \theutterance \stepcounter{utterance}  
    & & & \multicolumn{2}{p{0.3\linewidth}}{
        \cellcolor[rgb]{0.9,0.9,0.9}{
            \makecell[{{p{\linewidth}}}]{
                \texttt{\tiny{[GM$|$GM]}}
                \texttt{proposal} \\
            }
        }
    }
    & & \\ \\

    \theutterance \stepcounter{utterance}  
    & & & \multicolumn{2}{p{0.3\linewidth}}{
        \cellcolor[rgb]{0.9,0.9,0.9}{
            \makecell[{{p{\linewidth}}}]{
                \texttt{\tiny{[GM$|$GM]}}
                [1, 1, 3]
            }
        }
    }
    & & \\ \\

    \theutterance \stepcounter{utterance}  
    & & \multicolumn{4}{p{0.6\linewidth}}{
        \cellcolor[rgb]{0.9,0.9,0.9}{
            \makecell[{{p{\linewidth}}}]{
                \texttt{\tiny{[GM$\rangle$P2]}}
                \texttt{The other player has submitted a secret proposal. It is now your turn to submit a proposal of your own.} \\
            }
        }
    }
    & \\ \\

    \theutterance \stepcounter{utterance}  
    & & & \multicolumn{4}{p{0.6\linewidth}}{
        \cellcolor[rgb]{0.9,0.9,0.9}{
            \makecell[{{p{\linewidth}}}]{
                \texttt{\tiny{[GM$\langle$P2]}}
                \texttt{{[}Proposal: 3 canvas{]}} \\
            }
        }
    }
     \\ \\

    \theutterance \stepcounter{utterance}  
    & & & \multicolumn{2}{p{0.3\linewidth}}{
        \cellcolor[rgb]{0.9,0.9,0.9}{
            \makecell[{{p{\linewidth}}}]{
                \texttt{\tiny{[GM$|$GM]}}
                \texttt{proposal} \\
            }
        }
    }
    & & \\ \\

    \theutterance \stepcounter{utterance}  
    & & & \multicolumn{2}{p{0.3\linewidth}}{
        \cellcolor[rgb]{0.9,0.9,0.9}{
            \makecell[{{p{\linewidth}}}]{
                \texttt{\tiny{[GM$|$GM]}}
                [0, 0, 3]
            }
        }
    }
    & & \\ \\

    \theutterance \stepcounter{utterance}  
    & & & \multicolumn{2}{p{0.3\linewidth}}{
        \cellcolor[rgb]{0.9,0.9,0.9}{
            \makecell[{{p{\linewidth}}}]{
                \texttt{\tiny{[GM$|$GM]}}
                [[1, 1, 3], [0, 0, 3]]
            }
        }
    }
    & & \\ \\

    \theutterance \stepcounter{utterance}  
    & \multicolumn{4}{p{0.6\linewidth}}{
        \cellcolor[rgb]{0.95,0.95,0.95}{
            \makecell[{{p{\linewidth}}}]{
                \texttt{\tiny{[P2$\rightarrow$P1]}}
                \texttt{Post{-}episode Reflexion (not sent to GM):} \\
\texttt{New plan: I failed because I misinterpreted the user's message as a secret proposal when it was actually a public statement of their values and strategy. By responding with an analysis of their message and a counter{-}proposal in square brackets, I inadvertently triggered a conflict where both players claimed the same items (all canvases, candy, and glue), resulting in a non{-}complementary outcome and zero points. Next time, I will strictly distinguish between public communication and secret proposals. I will first send a public message acknowledging their stated values and proposing a specific split (e.g., "I propose you take the candy and glue, and I take the 3 canvases") without using square brackets. Only after receiving a clear, formatted secret proposal from them will I submit my own secret proposal in the exact required format `{[}Proposal: ...{]}` to ensure our claims are complementary.} \\
            }
        }
    }
    & & \\ \\

    \theutterance \stepcounter{utterance}  
    & & & \multicolumn{4}{p{0.6\linewidth}}{
        \cellcolor[rgb]{0.95,0.95,0.95}{
            \makecell[{{p{\linewidth}}}]{
                \texttt{\tiny{[P1$\rightarrow$P2]}}
                \texttt{Post{-}episode Reflexion (not sent to GM):} \\
\texttt{New plan: I failed because I sent my proposal inside square brackets (`{[}Proposal: 3 canvas{]}`) while the rules explicitly stated that square brackets are reserved *only* for secret proposals and that direct messages to the other player must not include them. By using brackets, the system likely interpreted my message as a secret proposal rather than a public negotiation message, causing the negotiation to fail or be misaligned. Next time, I will communicate my intent to take the canvases without any brackets (e.g., "I propose taking all 3 canvases") to ensure it is treated as a valid public message, and I will only use the exact format `{[}Proposal: ...{]}` if I decide to make a secret proposal immediately.} \\
            }
        }
    }
     \\ \\

\end{supertabular}
}

\end{document}
